\usepackage[T1]{fontenc}
\usepackage[utf8]{inputenc}
\usepackage{mathpazo}
\usepackage[scaled=.90]{helvet}
\usepackage{geometry,microtype,mathtools,amssymb}
\usepackage{booktabs,tabularx,array,multicol,longtable}
\usepackage{enumitem,xcolor,graphicx}
\usepackage[most]{tcolorbox}
\usepackage{tikz,pgfplots}
\usetikzlibrary{arrows.meta,positioning,calc,patterns,decorations.pathreplacing,trees}
\pgfplotsset{compat=1.18}
\usepackage{fancyhdr,lastpage,hyperref,titlesec}
\usepackage{listings}
\usepackage{needspace}

\geometry{inner=21mm,outer=17mm,top=20mm,bottom=22mm,headheight=15pt,footskip=10mm}
\definecolor{Navy}{HTML}{123B5D}
\definecolor{Blue}{HTML}{1D70A2}
\definecolor{Sky}{HTML}{EAF5FB}
\definecolor{Green}{HTML}{2E7D62}
\definecolor{Mint}{HTML}{EAF6F1}
\definecolor{Orange}{HTML}{C66528}
\definecolor{Cream}{HTML}{FFF4E8}
\definecolor{Red}{HTML}{A43E3E}
\definecolor{Rose}{HTML}{FBEDED}
\definecolor{Gray}{HTML}{F2F4F6}
\definecolor{Ink}{HTML}{1F2933}
\definecolor{Gold}{HTML}{D8A02B}
\definecolor{TeacherBlue}{HTML}{145DA0}
\color{Ink}
\hypersetup{colorlinks=true,linkcolor=Navy,urlcolor=Blue,pdftitle={Livret complet de mathématiques - Terminale STMG 2026-2027}}
\setlist[itemize]{leftmargin=1.45em,itemsep=2.5pt,topsep=4pt}
\setlist[enumerate]{leftmargin=1.85em,itemsep=4pt,topsep=4pt}
\setlength{\parindent}{0pt}
\setlength{\parskip}{4pt}
\raggedbottom
\frenchspacing
\renewcommand{\chaptername}{Chapitre}
\renewcommand{\contentsname}{Sommaire}
\renewcommand{\appendixname}{Annexe}
\titleformat{\chapter}[display]
  {\sffamily\color{Navy}}
  {\colorbox{Orange}{\color{white}\bfseries\Large\strut\ Chapitre \thechapter\ }}
  {10pt}{\bfseries\fontsize{27}{31}\selectfont}
  [\vspace{4pt}\color{Orange}\titlerule]
\titlespacing*{\chapter}{0pt}{-12pt}{24pt}
\titleformat{\section}{\sffamily\bfseries\Large\color{Navy}}{\thesection}{.7em}{}
\titleformat{\subsection}{\sffamily\bfseries\large\color{Blue}}{\thesubsection}{.7em}{}

\pagestyle{fancy}
\fancyhf{}
\fancyhead[LE,RO]{\small\sffamily\color{Navy}\thepage\,/\,\pageref*{LastPage}}
\fancyhead[LO]{\small\sffamily\color{Navy}\rightmark}
\fancyhead[RE]{\small\sffamily\color{Navy}\leftmark}
\renewcommand{\headrulewidth}{0.3pt}
\renewcommand{\chaptermark}[1]{\markboth{#1}{}}
\renewcommand{\sectionmark}[1]{\markright{#1}}

\newif\ifprof
\providecommand{\EditionProf}{false}
\edef\tempedition{\EditionProf}
\def\vrai{true}
\ifx\tempedition\vrai\proftrue\else\proffalse\fi

\tcbset{livretbox/.style={enhanced,breakable,arc=1.3mm,boxrule=.65pt,
  left=2.8mm,right=2.8mm,top=2.2mm,bottom=2.2mm,before skip=7pt,after skip=7pt,
  fonttitle=\bfseries\sffamily,coltitle=white,attach boxed title to top left={xshift=2mm,yshift=-2.2mm},
  boxed title style={boxrule=0pt,arc=.8mm,left=2mm,right=2mm,top=.8mm,bottom=.8mm}}}
\newtcolorbox{objectifs}{livretbox,colback=Sky,colframe=Blue,title={Objectifs de la séance},boxed title style={colback=Blue}}
\newtcolorbox{prerequis}{livretbox,colback=Gray,colframe=black!45,title={Rappel et prérequis},boxed title style={colback=black!60}}
\newtcolorbox{automatismes}{livretbox,colback=Cream,colframe=Orange,title={Automatismes -- 6 à 8 minutes, sans calculatrice},boxed title style={colback=Orange}}
\newtcolorbox{activite}[1][]{livretbox,colback=white,colframe=Green,title={Activité de recherche #1},boxed title style={colback=Green}}
\newtcolorbox{traceecrite}{livretbox,colback=Mint,colframe=Green,title={Trace écrite},boxed title style={colback=Green}}
\newtcolorbox{methode}{livretbox,colback=Sky,colframe=Navy,title={Méthode},boxed title style={colback=Navy}}
\newtcolorbox{exemple}{livretbox,colback=white,colframe=Blue,title={Exemple rédigé},boxed title style={colback=Blue}}
\newtcolorbox{exercice}[1][]{livretbox,colback=white,colframe=black!45,title={Entraînement #1},boxed title style={colback=black!65}}
\newtcolorbox{exerciceapp}{livretbox,colback=white,colframe=Gold,title={Exercice d'application},boxed title style={colback=Gold!85!black}}
\newtcolorbox{probleme}[1][]{livretbox,colback=white,colframe=Navy,title={Problème #1},boxed title style={colback=Navy}}
\newtcolorbox{manuel}{livretbox,colback=Gray,colframe=black!40,title={Repères dans le manuel},boxed title style={colback=black!60}}
\newtcolorbox{dm}{livretbox,colback=Cream,colframe=Orange,title={Travail personnel},boxed title style={colback=Orange}}
\newtcolorbox{sortie}{livretbox,colback=Rose,colframe=Red,title={Billet de sortie -- 3 minutes},boxed title style={colback=Red}}
\newtcolorbox{correction}{livretbox,colback=blue!3,colframe=TeacherBlue,title={Correction détaillée -- édition professeur},boxed title style={colback=TeacherBlue}}
\newtcolorbox{erreur}{enhanced,breakable,colback=Rose,colframe=Red,title={Erreur fréquente},fonttitle=\bfseries\sffamily,boxrule=.6pt,arc=2mm}
\newtcolorbox{retenir}{enhanced,breakable,colback=Mint,colframe=Green,title={À retenir},fonttitle=\bfseries\sffamily,boxrule=.8pt,arc=2mm}

\newcommand{\enseignant}[1]{\ifprof\begin{correction}#1\end{correction}\fi}
\newcommand{\rep}[1]{\ifprof\textcolor{Blue}{\textbf{#1}}\else\dotfill\fi}
\newcommand{\zoneeleve}[1]{\ifprof\else\par\vspace{2pt}\foreach \ligne in {1,...,#1}{\noindent\textcolor{black!22}{\rule{\linewidth}{.35pt}}\par\vspace{4.5mm}}\fi}
\newcommand{\difficulte}[1]{\hfill{\small\sffamily\color{Orange}#1}}
\newcommand{\competences}[1]{{\small\sffamily\textbf{Compétences :} #1}}
% Original application exercises (manuel/exo-A..D): \DefExo{session}{target}{statement}{teacher solution}{answer lines},
% then one \ExoSession{session} per session sheet. The target is shown to teachers only.
\makeatletter
\NewDocumentCommand{\DefExo}{m m +m +m m}{%
  \expandafter\long\expandafter\def\csname exo@#1\endcsname{%
    \begin{exerciceapp}
      #3
      \zoneeleve{#5}
    \end{exerciceapp}
    \enseignant{\textbf{Cible :} #2.\par #4}%
  }%
}
\newcommand{\ExoSession}[1]{%
  \@ifundefined{exo@#1}
    {\begin{erreur}Exercice à compléter pour la séance #1.\end{erreur}}
    {\csname exo@#1\endcsname}%
}
\makeatother

\newcommand{\manuelref}[1]{\begin{manuel}\textbf{Pour prolonger :} #1. L'exercice d'application ci-dessous est original ; ces références renvoient aux exercices du manuel pour s'entraîner davantage et retrouver ses figures.\end{manuel}}
\newcommand{\sessionheader}[5]{%
  \clearpage
  \phantomsection\addcontentsline{toc}{section}{Séance #1 -- #4}
  \begin{tcolorbox}[enhanced,colback=Navy,colframe=Navy,boxrule=0pt,arc=1.5mm,
    left=4mm,right=4mm,top=3.2mm,bottom=3.2mm]
  {\color{white}\sffamily\bfseries\Large Séance #1\hfill #4}\par\vspace{2mm}
  {\color{white!88}\sffamily\small #5\hfill #2\quad\textbullet\quad #3}
  \end{tcolorbox}}
\newcommand{\dsheader}[4]{%
  \clearpage\section*{Devoir surveillé #1 -- Version #2}
  \addcontentsline{toc}{section}{DS #1 -- Version #2}
  \begin{tcolorbox}[enhanced,colback=Navy,colframe=Navy,boxrule=0pt,arc=2mm]
  \color{white}\sffamily\textbf{#3}\hfill\textbf{Durée : 1 h}\quad\textbf{Note : /20}\\
  Nom : \rule{6cm}{.4pt}\hfill Classe : \rule{2.5cm}{.4pt}
  \end{tcolorbox}
  \competences{#4}\par\smallskip
  \textit{La calculatrice est autorisée sauf indication contraire. Toute démarche doit être lisible et les résultats contextualisés.}}
\newcommand{\points}[1]{\hfill{\small\sffamily\color{Navy}[#1 points]}}
\newcommand{\deuxcolonnes}[2]{\begin{minipage}[t]{.48\textwidth}#1\end{minipage}\hfill\begin{minipage}[t]{.48\textwidth}#2\end{minipage}}
\newcommand{\R}{\mathbb{R}}
\newcommand{\N}{\mathbb{N}}
\DeclareMathOperator{\e}{e}

\lstdefinestyle{python}{language=Python,basicstyle=\ttfamily\small,keywordstyle=\color{Blue}\bfseries,stringstyle=\color{Green},commentstyle=\color{black!55},showstringspaces=false,frame=single,rulecolor=\color{black!25},backgroundcolor=\color{Gray},breaklines=true}
